\RequirePackage{plautopatch}
\documentclass[platex,dvipdfmx,a4paper]{jsarticle}

\usepackage{amsmath}
\usepackage{amsfonts}
\usepackage{amssymb}

\usepackage[dvipdfmx]{graphicx}
\usepackage{float}

\usepackage{tikz}
\usepackage{caption}

% 「%」は以降の内容を「改行コードも含めて」無視するコマンド
\usepackage[%
 dvipdfmx,% 欧文ではコメントアウトする
 setpagesize=false,%
 bookmarks=true,%
 bookmarksdepth=tocdepth,%
 bookmarksnumbered=true,%
 colorlinks=false,%
 pdftitle={},%
 pdfsubject={},%
 pdfauthor={},%
 pdfkeywords={}%
]{hyperref}
% PDFのしおり機能の日本語文字化けを防ぐ((u)pLaTeXのときのみかく)
\usepackage{pxjahyper}

\newcommand{\hb}[2]{\frac{\partial #1}{\partial #2}}
\newcommand{\bbn}[2]{\frac{d #1}{d #2}}
\newcommand{\nhb}[3]{\frac{\partial^{#1} #2}{\partial #3^{#1}}}
\newcommand{\nbbn}[3]{\frac{d^{#1} #2}{d #3^{#1}}}
\newcommand{\vectxy}[2]{\begin{pmatrix}#1\\#2\end{pmatrix}}
\newcommand{\vectxyz}[3]{\begin{pmatrix}#1\\#2\\#3\end{pmatrix}}
\newcommand{\vect}[1]{\boldsymbol{\mathrm{#1}}}

\newcommand{\ssum}[3]{\sum_{#1=#2}^#3}

\newcommand{\veff}{V_{\mathrm{eff}}(r)}

\title{解析力学演習レポート1}
\author{嶋田 悠音}
\date{2026年6月16日}

\begin{document}
  \maketitle
  \section{問題1}
  $\nbbn{m}{u}{x}$

  \subsection*{(i)}
  本小問では、Legendreの微分方程式の解を$m$回微分して得られる関数が、Legendreの陪微分方程式を満たすことを示す。Legendreの微分方程式は次である：
  \begin{align*}
    (1-x^2)u''-2xu'+\nu(\nu+1)u=0。\tag{1}
  \end{align*}
  ここで、$v=\bbn[m]{u}{x}$とおく。(1)式を$x$で$m$回微分する。Leibnizの公式を用いると、
  \begin{align*}
    \bbn[m]{}{x}\{(1-x^2)u''\}=(1-x^2)v''-2mxv'-m(m-1)v
  \end{align*}
  であり、
  \begin{align*}
    \bbn[m]{}{x}\{-2xu'\}=-2xv'-2mv
  \end{align*}
  である。よって、$v$は次の方程式を満たす：
  \begin{align*}
    (1-x^2)v''-2(m+1)xv'+\{\nu(\nu+1)-m(m+1)\}v=0。
  \end{align*}
  次に、問題で定義された関数
  \begin{align*}
    y=(1-x^2)^{m/2}v
  \end{align*}
  が満たす方程式を求める。$s=1-x^2$とおくと、$v=s^{-m/2}y$である。これを微分する：
  \begin{align*}
    v'=s^{-m/2}y'+mxs^{-m/2-1}y
  \end{align*}
  であり、さらに微分して整理すると：
  \begin{align*}
    v''=s^{-m/2}y''+2mxs^{-m/2-1}y'+ms^{-m/2-1}y+m(m+2)x^2s^{-m/2-2}y。
  \end{align*}
  これらを$v$の方程式に代入する：
  \begin{align*}
    0&=(1-x^2)v''-2(m+1)xv'+\{\nu(\nu+1)-m(m+1)\}v \\
     &=s^{-m/2}\left[s y''-2xy'+\left\{\nu(\nu+1)-\frac{m^2}{s}\right\}y\right]。
  \end{align*}
  $s=1-x^2$であり、$s^{-m/2}\neq0$であるから、$y$は次の方程式を満たす：
  \begin{align*}
    (1-x^2)y''-2xy'+\left\{\nu(\nu+1)-\frac{m^2}{1-x^2}\right\}y=0。\tag{2}
  \end{align*}
  以上より、$y=(1-x^2)^{m/2}\bbn[m]{u}{x}$はLegendreの陪微分方程式を満たすことが示された。
  
  \subsection*{(ii)}
  本小問では、3次元Laplace方程式の一価連続な一般解を、球座標における変数分離によって求める。球座標におけるLaplace作用素は次である：
  \begin{align*}
    \Delta=\frac{1}{r^2}\bbn{}{r}\left(r^2\bbn{}{r}\right)+\frac{1}{r^2\sin\theta}\bbn{}{\theta}\left(\sin\theta\bbn{}{\theta}\right)+\frac{1}{r^2\sin^2\theta}\bbn[2]{}{\phi}。
  \end{align*}
  解を$u(r,\theta,\phi)=R(r)\Theta(\theta)\Phi(\phi)$とおき、Laplace方程式$\Delta u=0$に代入して$R\Theta\Phi$で割る。さらに両辺に$r^2$を掛けると：
  \begin{align*}
    \frac{1}{R}\bbn{}{r}\left(r^2\bbn{R}{r}\right)+\frac{1}{\Theta\sin\theta}\bbn{}{\theta}\left(\sin\theta\bbn{\Theta}{\theta}\right)+\frac{1}{\Phi\sin^2\theta}\bbn[2]{\Phi}{\phi}=0。
  \end{align*}
  まず$\phi$依存性を分離する。$\phi$方向の方程式を
  \begin{align*}
    \bbn[2]{\Phi}{\phi}+m^2\Phi=0
  \end{align*}
  とおくと、一価性$\Phi(\phi+2\pi)=\Phi(\phi)$より$m$は整数であり、解は$e^{im\phi}$で表される。したがって：
  \begin{align*}
    \Phi_m(\phi)=e^{im\phi},\qquad m\in\mathbb{Z}。
  \end{align*}
  残った$r,\theta$依存性について分離定数$l(l+1)$を導入すると、動径方程式と角度方程式は次のようになる：
  \begin{align*}
    \bbn{}{r}\left(r^2\bbn{R}{r}\right)-l(l+1)R=0
  \end{align*}
  および
  \begin{align*}
    \frac{1}{\sin\theta}\bbn{}{\theta}\left(\sin\theta\bbn{\Theta}{\theta}\right)+\left\{l(l+1)-\frac{m^2}{\sin^2\theta}\right\}\Theta=0。
  \end{align*}
  動径方程式に$R=r^\alpha$を代入する：
  \begin{align*}
    \alpha(\alpha+1)-l(l+1)=0。
  \end{align*}
  よって、動径方向の解は：
  \begin{align*}
    R(r)=A_{l,m}r^l+B_{l,m}r^{-(l+1)}
  \end{align*}
  である。次に角度方程式で$x=\cos\theta$とおくと、$\Theta(\theta)$はLegendreの陪微分方程式を満たす。よって、正則性を満たす角度関数は$P_l^{|m|}(\cos\theta)$で与えられる。また、$P_l^{|m|}$が多項式として正則であるためには$l=0,1,2,\ldots$かつ$|m|\leq l$である。したがって、3次元Laplace方程式の一価連続な一般解は：
  \begin{align*}
    u(r,\theta,\phi)=\ssum{l}{0}{\infty}\ssum{m}{-l}{l}\left(A_{l,m}r^l+B_{l,m}r^{-(l+1)}\right)P_l^{|m|}(\cos\theta)e^{im\phi}
  \end{align*}
  と書ける。
  
  \subsection*{(iii)}
  本小問では、半径1の球面上で$V=\cos\theta\cos\phi$に保たれているとき、球外の電位を求める。球外には電荷がないので、電位はLaplace方程式を満たす。また、無限遠で電位が発散しないため、球外解では$r^l$の項を除き、次の形を用いる：
  \begin{align*}
    V(r,\theta,\phi)=\ssum{l}{0}{\infty}\ssum{m}{-l}{l}B_{l,m}r^{-(l+1)}P_l^{|m|}(\cos\theta)e^{im\phi}。
  \end{align*}
  境界条件は$r=1$で
  \begin{align*}
    \ssum{l}{0}{\infty}\ssum{m}{-l}{l}B_{l,m}P_l^{|m|}(\cos\theta)e^{im\phi}=\cos\theta\cos\phi
  \end{align*}
  である。右辺は$\phi$について
  \begin{align*}
    \cos\theta\cos\phi=\frac{\cos\theta}{2}\left(e^{i\phi}+e^{-i\phi}\right)
  \end{align*}
  と書けるので、$m=\pm1$の成分のみが残る。そこで
  \begin{align*}
    \cos\theta=\ssum{l}{1}{\infty}a_lP_l^1(\cos\theta)
  \end{align*}
  と展開する。Legendre陪関数の直交性より、係数$a_l$は次のように求まる：
  \begin{align*}
    a_l=\frac{2l+1}{2}\frac{(l-1)!}{(l+1)!}\int_{-1}^{1}xP_l^1(x)\,dx。
  \end{align*}
  したがって、$B_{l,1}=B_{l,-1}=a_l/2$である。よって、球外の電位は次のように求まる：
  \begin{align*}
    V(r,\theta,\phi)=\ssum{l}{1}{\infty}a_l r^{-(l+1)}P_l^1(\cos\theta)\cos\phi,\qquad a_l=\frac{2l+1}{2}\frac{(l-1)!}{(l+1)!}\int_{-1}^{1}xP_l^1(x)\,dx。
  \end{align*}
  これが、球面上の境界条件$V=\cos\theta\cos\phi$に対応する球外電位である。
  
  \section{問題2}
  
  \subsection*{(i)}
  本小問では、Rodriguesの公式からLegendre多項式の複素積分表示を導く。Rodriguesの公式は：
  \begin{align*}
    P_n(z)=\frac{1}{2^n n!}\bbn[n]{}{z}(z^2-1)^n
  \end{align*}
  である。ここで$f(w)=(w^2-1)^n$とおく。Cauchyの積分公式より、$z$の周りを反時計回りに一周する閉曲線$C$に対して：
  \begin{align*}
    f^{(n)}(z)=\frac{n!}{2\pi i}\int_C\frac{f(w)}{(w-z)^{n+1}}\,dw
  \end{align*}
  が成り立つ。これをRodriguesの公式に代入する：
  \begin{align*}
    P_n(z)=\frac{1}{2^n n!}\frac{n!}{2\pi i}\int_C\frac{(w^2-1)^n}{(w-z)^{n+1}}\,dw=\frac{1}{2\pi i}\int_C\frac{(w^2-1)^n}{2^n(w-z)^{n+1}}\,dw。
  \end{align*}
  以上より、Legendre多項式の複素積分表示が示された。
  
  \subsection*{(ii)}
  本小問では、Legendre多項式の母関数を導く。(i)で得た積分表示を用いる：
  \begin{align*}
    \ssum{n}{0}{\infty}P_n(z)\xi^n=\frac{1}{2\pi i}\int_C\ssum{n}{0}{\infty}\frac{(w^2-1)^n\xi^n}{2^n(w-z)^{n+1}}\,dw。
  \end{align*}
  $\xi$が十分小さいとして、和と積分を交換する。右辺の級数は等比級数であり：
  \begin{align*}
    \ssum{n}{0}{\infty}\frac{(w^2-1)^n\xi^n}{2^n(w-z)^{n+1}}=\frac{1}{w-z-\frac{\xi}{2}(w^2-1)}
  \end{align*}
  である。したがって、
  \begin{align*}
    \ssum{n}{0}{\infty}P_n(z)\xi^n=\frac{1}{2\pi i}\int_C\frac{dw}{w-z-\frac{\xi}{2}(w^2-1)}。
  \end{align*}
  分母の零点を調べる：
  \begin{align*}
    w-z-\frac{\xi}{2}(w^2-1)=0。
  \end{align*}
  この方程式の解のうち、$\xi\to0$で$w=z$に近づく根だけが、積分路$C$の内部にある。この根を$w_0$とする。分母を$F(w)$と書くと、留数定理より：
  \begin{align*}
    \ssum{n}{0}{\infty}P_n(z)\xi^n=\frac{1}{F'(w_0)}。
  \end{align*}
  ここで$F'(w)=1-\xi w$である。また、根の公式から
  \begin{align*}
    w_0=\frac{1-\sqrt{1-2z\xi+\xi^2}}{\xi}
  \end{align*}
  であり、したがって：
  \begin{align*}
    F'(w_0)=1-\xi w_0=\sqrt{1-2z\xi+\xi^2}。
  \end{align*}
  よって、
  \begin{align*}
    \frac{1}{\sqrt{1-2z\xi+\xi^2}}=\ssum{n}{0}{\infty}P_n(z)\xi^n
  \end{align*}
  が得られた。これがLegendre多項式の母関数である。
  
  \subsection*{(iii)}
  本小問では、$z$軸上の点$z=a$に置かれた点電荷$q$が作る電位をLegendre多項式で展開する。点電荷から点$\vect{r}$までの距離は：
  \begin{align*}
    |\vect{r}-a\vect{e}_z|=\sqrt{r^2+a^2-2ar\cos\theta}
  \end{align*}
  である。したがって、電位は：
  \begin{align*}
    V(\vect{r})=\frac{q}{4\pi\epsilon_0}\frac{1}{\sqrt{r^2+a^2-2ar\cos\theta}}
  \end{align*}
  である。ここで、
  \begin{align*}
    r_>=\max(r,a),\qquad r_<=\min(r,a)
  \end{align*}
  とおくと、分母は次のように書ける：
  \begin{align*}
    \sqrt{r^2+a^2-2ar\cos\theta}=r_>\sqrt{1-2\cos\theta\frac{r_<}{r_>}+\left(\frac{r_<}{r_>}\right)^2}。
  \end{align*}
  問題2(ii)の母関数を$z=\cos\theta,\ \xi=r_</r_>$として用いる：
  \begin{align*}
    \frac{1}{\sqrt{1-2\cos\theta\frac{r_<}{r_>}+\left(\frac{r_<}{r_>}\right)^2}}=\ssum{n}{0}{\infty}\left(\frac{r_<}{r_>}\right)^nP_n(\cos\theta)。
  \end{align*}
  したがって、点電荷の電位は：
  \begin{align*}
    V(\vect{r})=\frac{q}{4\pi\epsilon_0 r_>}\ssum{n}{0}{\infty}\left(\frac{r_<}{r_>}\right)^nP_n(\cos\theta)。
  \end{align*}
  ただし、
  \begin{align*}
    r_>=r,\ r_<=a\quad(r>a),\qquad r_>=a,\ r_<=r\quad(r<a)
  \end{align*}
  である。
  
  \subsection*{(iv)}
  本小問では、点$z=a$に電荷$+q$、点$z=-a$に電荷$-q$を置いたとき、遠方$r\gg a$での電位を求める。まず、点$z=a$にある電荷$+q$による電位は：
  \begin{align*}
    V_+(\vect{r})=\frac{q}{4\pi\epsilon_0 r}\ssum{n}{0}{\infty}\left(\frac{a}{r}\right)^nP_n(\cos\theta)
  \end{align*}
  である。一方、点$z=-a$にある電荷$-q$による電位は、$\cos\theta$を$-\cos\theta$に置き換えればよいので：
  \begin{align*}
    V_-(\vect{r})=-\frac{q}{4\pi\epsilon_0 r}\ssum{n}{0}{\infty}\left(\frac{a}{r}\right)^nP_n(-\cos\theta)
  \end{align*}
  である。Legendre多項式の性質$P_n(-x)=(-1)^nP_n(x)$を用いると、全電位は：
  \begin{align*}
    V(\vect{r})=\frac{q}{4\pi\epsilon_0 r}\ssum{n}{0}{\infty}\left(\frac{a}{r}\right)^n\{1-(-1)^n\}P_n(\cos\theta)。
  \end{align*}
  偶数$n$の項は消え、奇数$n$の項だけが残る。遠方$r\gg a$では最低次の$n=1$の項が支配的であるから：
  \begin{align*}
    V(\vect{r})=\frac{q}{4\pi\epsilon_0 r}\,2\frac{a}{r}P_1(\cos\theta)+O\left(\frac{a^3}{r^4}\right)=\frac{1}{4\pi\epsilon_0}\frac{2qa\cos\theta}{r^2}+O\left(\frac{a^3}{r^4}\right)。
  \end{align*}
  したがって、遠方では双極子電位が得られる。
  
  \subsection*{おまけ}
  遠方で電位の角度依存性が$P_2(\cos\theta)$となるようにするには、全電荷と双極子モーメントを消し、四重極成分を最初に残せばよい。例えば、$z=a$に電荷$+q$、$z=-a$に電荷$+q$、原点に電荷$-2q$を置く配置を考える。このとき遠方の電位は：
  \begin{align*}
    V(\vect{r})=\frac{q}{4\pi\epsilon_0 r}\ssum{n}{0}{\infty}\left(\frac{a}{r}\right)^n\{1+(-1)^n\}P_n(\cos\theta)-\frac{2q}{4\pi\epsilon_0 r}。
  \end{align*}
  $n=0$成分は原点の電荷$-2q$により打ち消され、奇数$n$成分は対称性により消える。したがって、遠方での最低次項は$n=2$であり：
  \begin{align*}
    V(\vect{r})=\frac{1}{4\pi\epsilon_0}\frac{2qa^2}{r^3}P_2(\cos\theta)+O\left(\frac{a^4}{r^5}\right)。
  \end{align*}
  よって、この配置により遠方での角度依存性を$P_2(\cos\theta)$にできる。
  
  \section{問題3}
  
  \subsection*{I(i)}
  本小問では、定義された球面調和関数について、$l=0,1$の場合を具体的に書き下す。球面調和関数は：
  \begin{align*}
    Y_{l,m}(\theta,\phi)=i^{|m|+m}\sqrt{\frac{2l+1}{4\pi}\frac{(l-|m|)!}{(l+|m|)!}}P_l^{|m|}(\cos\theta)e^{im\phi}
  \end{align*}
  で定義される。まず$l=0$では$m=0$のみであり、$P_0^0(\cos\theta)=1$であるから：
  \begin{align*}
    Y_{0,0}(\theta,\phi)=\frac{1}{\sqrt{4\pi}}。
  \end{align*}
  次に$l=1$の場合を考える。このとき$m=-1,0,1$である。$P_1^0(\cos\theta)=\cos\theta$、$P_1^1(\cos\theta)=\sin\theta$を用いる。まず$m=0$について：
  \begin{align*}
    Y_{1,0}(\theta,\phi)=\sqrt{\frac{3}{4\pi}}\cos\theta。
  \end{align*}
  次に$m=1$について、$i^{|1|+1}=i^2=-1$であるから：
  \begin{align*}
    Y_{1,1}(\theta,\phi)=-\sqrt{\frac{3}{8\pi}}\sin\theta\,e^{i\phi}。
  \end{align*}
  また$m=-1$について、$i^{|-1|-1}=i^0=1$であるから：
  \begin{align*}
    Y_{1,-1}(\theta,\phi)=\sqrt{\frac{3}{8\pi}}\sin\theta\,e^{-i\phi}。
  \end{align*}
  以上より、$l=0,1$の球面調和関数はすべて書き下された。
  
  \subsection*{I(ii)}
  本小問では、球面調和関数の加法定理を用いて、任意の点に置かれた点電荷の電位を展開する。問題2(iii)より、相対角$\gamma$を用いると点電荷の電位は：
  \begin{align*}
    V(\vect{r})=\frac{q}{4\pi\epsilon_0 r_>}\ssum{l}{0}{\infty}\left(\frac{r_<}{r_>}\right)^lP_l(\cos\gamma)
  \end{align*}
  と書ける。ここで球面調和関数の加法定理：
  \begin{align*}
    P_l(\cos\gamma)=\frac{4\pi}{2l+1}\ssum{m}{-l}{l}Y_{l,m}^*(\theta',\phi')Y_{l,m}(\theta,\phi)
  \end{align*}
  を用いる。この式を電位の展開に代入する：
  \begin{align*}
    V(\vect{r})=\frac{q}{4\pi\epsilon_0 r_>}\ssum{l}{0}{\infty}\left(\frac{r_<}{r_>}\right)^l\frac{4\pi}{2l+1}\ssum{m}{-l}{l}Y_{l,m}^*(\theta',\phi')Y_{l,m}(\theta,\phi)。
  \end{align*}
  したがって：
  \begin{align*}
    V(\vect{r})=\frac{q}{\epsilon_0 r_>}\ssum{l}{0}{\infty}\ssum{m}{-l}{l}\frac{1}{2l+1}\left(\frac{r_<}{r_>}\right)^lY_{l,m}^*(\theta',\phi')Y_{l,m}(\theta,\phi)
  \end{align*}
  を得る。
  
  \subsection*{II(iii)}
  本小問では、半径$a$の球面上に軸対称な電荷分布$q\cos\theta\,d\Omega$が置かれたとき、球の外側の電位を求める。ここで$d\Omega=\sin\theta\,d\theta\,d\phi$である。前小問の結果を連続分布に対して積分する。球外では$r>a$であるから、$r_>=r,\ r_<=a$であり：
  \begin{align*}
    V(r,\theta,\phi)=\frac{1}{\epsilon_0 r}\ssum{l}{0}{\infty}\ssum{m}{-l}{l}\frac{1}{2l+1}\left(\frac{a}{r}\right)^lY_{l,m}(\theta,\phi)\int q\cos\theta'Y_{l,m}^*(\theta',\phi')\,d\Omega'
  \end{align*}
  である。電荷分布は$\phi'$に依存しないので、$m=0$の項だけが残る。また、$\cos\theta'$は$Y_{1,0}$に比例するので、$l=1$のみが残る。実際、
  \begin{align*}
    Y_{1,0}(\theta,\phi)=\sqrt{\frac{3}{4\pi}}\cos\theta
  \end{align*}
  であるから、必要な積分は：
  \begin{align*}
    \int q\cos\theta'Y_{1,0}^*(\theta',\phi')\,d\Omega'=q\sqrt{\frac{3}{4\pi}}\int\cos^2\theta'\,d\Omega'=q\sqrt{\frac{3}{4\pi}}\frac{4\pi}{3}
  \end{align*}
  である。したがって、球外電位は：
  \begin{align*}
    V(r,\theta,\phi)=\frac{1}{\epsilon_0 r}\frac{1}{3}\frac{a}{r}\sqrt{\frac{3}{4\pi}}\cos\theta\cdot q\sqrt{\frac{3}{4\pi}}\frac{4\pi}{3}=\frac{qa}{3\epsilon_0}\frac{\cos\theta}{r^2}。
  \end{align*}
  以上より、球の外側での電位は：
  \begin{align*}
    V(r,\theta,\phi)=\frac{qa\cos\theta}{3\epsilon_0 r^2}\qquad (r>a)
  \end{align*}
  である。
  
  \subsection*{II(iv)}
  本小問では、水素原子の$1s$軌道にある電子の球対称な電荷分布が作る電位を求める。与えられた電荷分布は球対称であるから、電位は$r$のみに依存する。体積電荷密度を
  \begin{align*}
    \rho(r)=\frac{e}{\pi a_0^3}e^{-2r/a_0}
  \end{align*}
  と読むと、半径$r$より内側の電荷は：
  \begin{align*}
    Q_{\mathrm{in}}(r)=4\pi\int_0^r\rho(r')r'^2\,dr'=e\left\{1-e^{-2r/a_0}\left(1+\frac{2r}{a_0}+\frac{2r^2}{a_0^2}\right)\right\}
  \end{align*}
  である。また、半径$r$より外側の球殻が原点から距離$r$の点に作る電位への寄与は：
  \begin{align*}
    4\pi\int_r^\infty \rho(r')r'\,dr'=e\,e^{-2r/a_0}\left(\frac{2r}{a_0^2}+\frac{1}{a_0}\right)
  \end{align*}
  である。したがって、球対称分布の電位は：
  \begin{align*}
    V(r)=\frac{1}{4\pi\epsilon_0}\left\{\frac{Q_{\mathrm{in}}(r)}{r}+4\pi\int_r^\infty \rho(r')r'\,dr'\right\}
  \end{align*}
  である。上で求めた式を代入して整理する：
  \begin{align*}
    V(r)&=\frac{e}{4\pi\epsilon_0}\left[\frac{1}{r}\left\{1-e^{-2r/a_0}\left(1+\frac{2r}{a_0}+\frac{2r^2}{a_0^2}\right)\right\}+e^{-2r/a_0}\left(\frac{2r}{a_0^2}+\frac{1}{a_0}\right)\right] \\
        &=\frac{e}{4\pi\epsilon_0}\left\{\frac{1}{r}-e^{-2r/a_0}\left(\frac{1}{r}+\frac{1}{a_0}\right)\right\}。
  \end{align*}
  よって、水素原子の$1s$軌道の電子分布が作る電位は：
  \begin{align*}
    V(r)=\frac{e}{4\pi\epsilon_0 r}\left\{1-\left(1+\frac{r}{a_0}\right)e^{-2r/a_0}\right\}
  \end{align*}
  である。電子の電荷を$-e$として扱う場合は、右辺全体に負号を付ければよい。  
\end{document}